\section{Gap Analysis}
Despite advancements in signaling, several shortcomings remain in traditional interlocking systems. High costs make modern vendor solutions inaccessible for many developing countries, leaving them dependent on outdated technologies. Relay and PLC-based systems also lack flexibility; adapting them to new layouts or operational changes often requires extensive rewiring, which limits scalability.

Another persistent issue is the absence of effective simulation tools. Without digital sandboxes, safety logic is difficult to test and validate before deployment, increasing risks during commissioning and restricting opportunities for operator training. Legacy systems are further constrained by poor integration with modern data-driven platforms, preventing their use in predictive maintenance or centralized control strategies.

Perhaps most critically, interlocking logic in conventional systems is opaque, often hidden within wiring or PLC blocks, making it difficult to audit, extend, or regulate. These limitations highlight the need for a modular, transparent, and software-defined approach that combines traditional safety robustness with the agility required in modern railway operations.
